\chapter{Projects with HobbySpark}

This chapter puts the classes and hardware objects introduced in the previous
chapter to practical use. Each project combines one or more HobbySpark
classes to create a useful hardware application.

The projects are arranged so that the reader can gradually become comfortable
with combining different hardware objects. Each project contains an
introduction, the complete HobbySpark program, and a line-by-line explanation
of the code.

\section{Button-Controlled Buzzer}

\subsection{Introduction}

This project creates a simple button-controlled buzzer. When the button is
pressed, the active buzzer is turned on. When the button is released, the
buzzer is turned off.

This is a useful example because it demonstrates how a \texttt{Button} can
be combined with an \texttt{ActiveBuzzer}. It also introduces the use of
\texttt{when\_read()} to make the program respond automatically to a change
in the button's state.

\subsection{Code}

\begin{lstlisting}
from stub import *

set_board(ArduinoNano())

buzzer = ActiveBuzzer(4)
button = Button(2)

def button_pressed():
    buzzer.on()

button.when_read(button_pressed)
\end{lstlisting}

\subsection{Explanation}

\begin{enumerate}
    \item \texttt{from stub import *} imports the HobbySpark classes and
    functions required by the program.

    \item \texttt{set\_board(ArduinoNano())} tells HobbySpark that the program
    is intended to run on an Arduino Nano.

    \item \texttt{buzzer = ActiveBuzzer(4)} creates an
    \texttt{ActiveBuzzer} connected to pin 4.

    \item \texttt{button = Button(2)} creates a \texttt{Button} connected to
    pin 2.

    \item \texttt{def button\_pressed():} defines a function that will be
    executed when the button reading becomes true.

    \item \texttt{buzzer.on()} turns the active buzzer on.

    \item \texttt{button.when\_read(button\_pressed)} tells HobbySpark to run
    \texttt{button\_pressed} whenever the button reading is true.
\end{enumerate}


\newpage

\section{Electronic Doorbell}

\subsection{Introduction}

This project turns a button and an active buzzer into a simple electronic
doorbell. Pressing the button causes the buzzer to sound for one second.

Unlike the previous project, the buzzer is not left on indefinitely. The
program performs an action, waits for a fixed amount of time, and then turns
the buzzer off.

\subsection{Code}

\begin{lstlisting}
from stub import *

set_board(ArduinoNano())

buzzer = ActiveBuzzer(4)
button = Button(2)

def ring_bell():
    buzzer.on()
    wait(1)
    buzzer.off()

button.when_read(ring_bell)
\end{lstlisting}

\subsection{Explanation}

\begin{enumerate}
    \item \texttt{from stub import *} imports the HobbySpark framework.

    \item \texttt{set\_board(ArduinoNano())} selects the Arduino Nano as the
    target board.

    \item \texttt{buzzer = ActiveBuzzer(4)} creates an active buzzer connected
    to pin 4.

    \item \texttt{button = Button(2)} creates a button connected to pin 2.

    \item \texttt{def ring\_bell():} defines the function that represents the
    doorbell action.

    \item \texttt{buzzer.on()} switches the buzzer on.

    \item \texttt{wait(1)} waits for one second.

    \item \texttt{buzzer.off()} switches the buzzer off after the one-second
    delay.

    \item \texttt{button.when\_read(ring\_bell)} connects the button reading
    to the doorbell function. When the button reading is true,
    \texttt{ring\_bell()} is executed.
\end{enumerate}


\newpage

\section{Parking Distance Warning System}

\subsection{Introduction}

This project creates a simple parking-distance warning system using an
ultrasonic sensor, an RGB LED, and a passive buzzer.

The ultrasonic sensor measures the distance to an object. The RGB LED
indicates the state of the system, while the passive buzzer produces an
audible warning.

This is a practical example of combining several different HobbySpark
hardware objects into one application.

\subsection{Code}

\begin{lstlisting}
from stub import *

set_board(ArduinoNano())

sensor = UltrasonicSensor(6, 7)
led = RGBLed(3, 5, 9)
buzzer = PassiveBuzzer(10)

distance = sensor.distance()

if distance < 10:
    led.set(100, 0, 0)
    buzzer.write(1000, 1, SEC)
elif distance < 30:
    led.set(100, 100, 0)
    buzzer.write(600, 1, SEC)
else:
    led.set(0, 100, 0)
\end{lstlisting}

\subsection{Explanation}

\begin{enumerate}
    \item \texttt{from stub import *} imports the HobbySpark framework.

    \item \texttt{set\_board(ArduinoNano())} selects the Arduino Nano.

    \item \texttt{sensor = UltrasonicSensor(6, 7)} creates an ultrasonic
    sensor. Pin 6 is used as the trigger pin and pin 7 as the echo pin.

    \item \texttt{led = RGBLed(3, 5, 9)} creates an RGB LED using pins 3,
    5, and 9 for its red, green, and blue channels.

    \item \texttt{buzzer = PassiveBuzzer(10)} creates a passive buzzer on
    pin 10.

    \item \texttt{distance = sensor.distance()} obtains the measured distance
    from the ultrasonic sensor.

    \item \texttt{if distance < 10:} checks whether an object is less than
    10 centimeters away.

    \item \texttt{led.set(100, 0, 0)} sets the RGB LED to red when the object
    is very close.

    \item \texttt{buzzer.write(1000, 1, SEC)} plays a high-frequency warning
    tone for one second.

    \item \texttt{elif distance < 30:} checks whether the object is less than
    30 centimeters away but was not already found to be less than
    10 centimeters away.

    \item \texttt{led.set(100, 100, 0)} sets the RGB LED to yellow.

    \item \texttt{buzzer.write(600, 1, SEC)} produces a lower-frequency
    warning tone.

    \item \texttt{else:} handles objects that are at least 30 centimeters
    away.

    \item \texttt{led.set(0, 100, 0)} sets the RGB LED to green, indicating
    that the object is at a safer distance.
\end{enumerate}


\newpage

\section{Automatic Barrier Gate}

\subsection{Introduction}

This project creates a small automatic barrier gate. A button is used to
request the gate to open, and a servo motor moves the barrier to an open
position. After a short delay, the servo returns the barrier to its closed
position.

An RGB LED is also used to indicate the state of the gate.

\subsection{Code}

\begin{lstlisting}
from stub import *

set_board(ArduinoNano())

button = Button(2)
servo = ServoMotor(9)
led = RGBLed(3, 5, 6)

def open_gate():
    led.set(0, 100, 0)
    servo.write(90)
    wait(3)
    servo.write(0)
    led.off()

button.when_read(open_gate)
\end{lstlisting}

\subsection{Explanation}

\begin{enumerate}
    \item \texttt{from stub import *} imports the HobbySpark framework.

    \item \texttt{set\_board(ArduinoNano())} selects the Arduino Nano.

    \item \texttt{button = Button(2)} creates a button connected to pin 2.

    \item \texttt{servo = ServoMotor(9)} creates a servo motor connected to
    pin 9.

    \item \texttt{led = RGBLed(3, 5, 6)} creates an RGB LED using pins 3,
    5, and 6.

    \item \texttt{def open\_gate():} defines the operation performed when the
    button is pressed.

    \item \texttt{led.set(0, 100, 0)} changes the RGB LED to green to indicate
    that the gate is opening.

    \item \texttt{servo.write(90)} moves the servo to 90 degrees.

    \item \texttt{wait(3)} keeps the gate open for three seconds.

    \item \texttt{servo.write(0)} returns the servo to zero degrees,
    closing the gate.

    \item \texttt{led.off()} switches the RGB LED off.

    \item \texttt{button.when\_read(open\_gate)} causes the gate-opening
    function to execute when the button reading is true.
\end{enumerate}


\newpage

\section{Ultrasonic Security Alarm}

\subsection{Introduction}

This project demonstrates a simple security alarm. An ultrasonic sensor
monitors the distance in front of the device. When an object reaches the
specified distance, the program activates an active buzzer.

This example uses the \texttt{when\_distance()} method, allowing the
ultrasonic sensor to trigger a function when the specified distance
condition is reached.

\subsection{Code}

\begin{lstlisting}
from stub import *

set_board(ArduinoNano())

sensor = UltrasonicSensor(6, 7)
buzzer = ActiveBuzzer(4)
led = Led(8)

def alarm():
    buzzer.on()
    led.on()

sensor.when_distance(10, alarm)
\end{lstlisting}

\subsection{Explanation}

\begin{enumerate}
    \item \texttt{from stub import *} imports the HobbySpark framework.

    \item \texttt{set\_board(ArduinoNano())} selects the Arduino Nano.

    \item \texttt{sensor = UltrasonicSensor(6, 7)} creates an ultrasonic
    sensor using pins 6 and 7.

    \item \texttt{buzzer = ActiveBuzzer(4)} creates an active buzzer on
    pin 4.

    \item \texttt{led = Led(8)} creates an LED connected to pin 8.

    \item \texttt{def alarm():} defines the function that will be executed
    when the alarm condition occurs.

    \item \texttt{buzzer.on()} activates the buzzer.

    \item \texttt{led.on()} switches the warning LED on.

    \item \texttt{sensor.when\_distance(10, alarm)} tells HobbySpark to
    execute \texttt{alarm()} when the measured distance matches the
    specified distance.
\end{enumerate}


\newpage

\section{Serial-Controlled LCD Message System}

\subsection{Introduction}

This project combines the serial monitor with an LCD display. It demonstrates
how information can be read from the serial monitor and displayed on an
I2C LCD.

It is useful when a program needs to communicate with a user through both
the computer and a physical display.

\subsection{Code}

\begin{lstlisting}
from stub import *

set_board(ArduinoNano())

serial = SerialMoniter(9600)
display = I2C_LCD_Display(21, 22, 16, 2)

serial.print("Enter a message:")
message = serial.input()

display.set_cursor(0, 0)
display.write(message)
\end{lstlisting}

\subsection{Explanation}

\begin{enumerate}
    \item \texttt{from stub import *} imports the HobbySpark framework.

    \item \texttt{set\_board(ArduinoNano())} selects the Arduino Nano.

    \item \texttt{serial = SerialMoniter(9600)} creates a serial monitor
    object using a baud rate of 9600.

    \item \texttt{display = I2C\_LCD\_Display(21, 22, 16, 2)} creates a
    16-column by 2-row I2C LCD display. The SDA and SCL pins are 21 and 22.

    \item \texttt{serial.print("Enter a message:")} displays a prompt in
    the serial monitor.

    \item \texttt{message = serial.input()} reads the user's serial input
    and stores it in \texttt{message}.

    \item \texttt{display.set\_cursor(0, 0)} moves the LCD cursor to the
    first column of the first row.

    \item \texttt{display.write(message)} writes the received message to
    the LCD.
\end{enumerate}


\newpage

\section{RGB Status Indicator}

\subsection{Introduction}

This project demonstrates how an RGB LED can be used as a status indicator.
The program represents three different states using three different colors.

The project also demonstrates the \texttt{RGB} class. Instead of directly
passing three color components to the RGB LED, an \texttt{RGB} object is
used to represent a color.

\subsection{Code}

\begin{lstlisting}
from stub import *

set_board(ArduinoNano())

led = RGBLed(3, 5, 6)

ready = RGB(0, 100, 0)
warning = RGB(100, 100, 0)
error = RGB(100, 0, 0)

led.set(ready.r, ready.g, ready.b)
wait(1)

led.set(warning.r, warning.g, warning.b)
wait(1)

led.set(error.r, error.g, error.b)
\end{lstlisting}

\subsection{Explanation}

\begin{enumerate}
    \item \texttt{from stub import *} imports the HobbySpark framework.

    \item \texttt{set\_board(ArduinoNano())} selects the Arduino Nano.

    \item \texttt{led = RGBLed(3, 5, 6)} creates an RGB LED using three
    output pins.

    \item \texttt{ready = RGB(0, 100, 0)} creates an RGB color representing
    green.

    \item \texttt{warning = RGB(100, 100, 0)} creates an RGB color
    representing yellow.

    \item \texttt{error = RGB(100, 0, 0)} creates an RGB color representing
    red.

    \item \texttt{led.set(ready.r, ready.g, ready.b)} sets the LED to the
    ready color.

    \item \texttt{wait(1)} waits for one second.

    \item \texttt{led.set(warning.r, warning.g, warning.b)} changes the LED
    to the warning color.

    \item \texttt{wait(1)} waits for another second.

    \item \texttt{led.set(error.r, error.g, error.b)} changes the LED to the
    error color.
\end{enumerate}


\newpage

\section{Ultrasonic Distance Display}

\subsection{Introduction}

This project combines an ultrasonic sensor and an I2C LCD display to create
a simple distance display. The measured distance is written to the LCD.

This is a useful example of combining a sensor with a display rather than
simply printing sensor data to the terminal.

\subsection{Code}

\begin{lstlisting}
from stub import *

set_board(ArduinoNano())

sensor = UltrasonicSensor(6, 7)
display = I2C_LCD_Display(21, 22, 16, 2)

distance = sensor.distance()

display.set_cursor(0, 0)
display.write("Distance:")

display.set_cursor(0, 1)
display.write(distance)
\end{lstlisting}

\subsection{Explanation}

\begin{enumerate}
    \item \texttt{from stub import *} imports the HobbySpark framework.

    \item \texttt{set\_board(ArduinoNano())} selects the Arduino Nano.

    \item \texttt{sensor = UltrasonicSensor(6, 7)} creates an ultrasonic
    sensor using pins 6 and 7.

    \item \texttt{display = I2C\_LCD\_Display(21, 22, 16, 2)} creates the
    16 by 2 LCD display.

    \item \texttt{distance = sensor.distance()} obtains the measured distance.

    \item \texttt{display.set\_cursor(0, 0)} places the cursor at the
    beginning of the first row.

    \item \texttt{display.write("Distance:")} writes the label
    \texttt{Distance:} to the display.

    \item \texttt{display.set\_cursor(0, 1)} moves the cursor to the
    beginning of the second row.

    \item \texttt{display.write(distance)} writes the measured distance
    to the LCD.
\end{enumerate}


\newpage

\section{Servo Position Controller}

\subsection{Introduction}

This project uses a button to move a servo motor between two positions.
Each time the button is read as true, the servo is moved to a predefined
position.

The project demonstrates how a physical input can control a mechanical
output.

\subsection{Code}

\begin{lstlisting}
from stub import *

set_board(ArduinoNano())

button = Button(2)
servo = ServoMotor(9)

def move_servo():
    servo.write(90)

button.when_read(move_servo)
\end{lstlisting}

\subsection{Explanation}

\begin{enumerate}
    \item \texttt{from stub import *} imports the HobbySpark framework.

    \item \texttt{set\_board(ArduinoNano())} selects the Arduino Nano.

    \item \texttt{button = Button(2)} creates a button connected to pin 2.

    \item \texttt{servo = ServoMotor(9)} creates a servo motor connected to
    pin 9.

    \item \texttt{def move\_servo():} defines the function that moves the
    servo.

    \item \texttt{servo.write(90)} commands the servo to move to 90 degrees.

    \item \texttt{button.when\_read(move\_servo)} connects the button to
    the servo-control function.
\end{enumerate}
